% ============================================================
%  Глава 7 — ЛКР через DARE
% ============================================================
\chapter{ЛКР: оптимальный регулятор}
\label{ch:lqr}

\section{Вопрос}

В прошлой главе мы сказали: ЛКР даёт $\vect{u}_k = -\mat{K}\vect{x}_k$,
и $\mat{K}$ считается через \emph{уравнение Риккати}. Сейчас аккуратно
разберём, \emph{что} это за уравнение и \emph{откуда} оно берётся.

\begin{ideabox}
ЛКР --- это \emph{задача оптимизации} на бесконечном горизонте.
Решить её «в лоб» (например, методом Беллмана) можно, но получается
тяжёлая рекурсия. Чудо в том, что \textbf{для линейного объекта с
квадратичным критерием эта рекурсия сходится к одной матрице $\mat{P}$,
которую можно найти из алгебраического уравнения}.

Это уравнение --- DARE (Discrete-time Algebraic Riccati Equation).
\end{ideabox}

\section{Постановка задачи}

Объект:
\[
  \vect{x}_{k+1} = \mat{A}_d\,\vect{x}_k + \mat{B}_d\,\vect{u}_k,
  \qquad \vect{x}_0 \text{ задано}.
\]

Критерий:
\[
  J = \sum_{k=0}^{\infty} \Big(\vect{x}_k\T \mat{Q}\,\vect{x}_k +
                                  \vect{u}_k\T \mat{R}\,\vect{u}_k\Big).
\]

Цель: найти последовательность $\{\vect{u}_k\}_{k=0}^\infty$, которая
\textbf{минимизирует} $J$.

\section{Принцип динамического программирования}

\begin{ideabox}
\textbf{Принцип Беллмана.} Оптимальная политика на бесконечном
горизонте обладает свойством: \emph{какова бы ни была начальная
ситуация и первое решение, последующие решения должны быть
оптимальными для оставшейся задачи}.

Для квадратичного критерия и линейной динамики из этого принципа
следует, что \textbf{функция ценности} (стоимость до бесконечности
из состояния $\vect{x}$) тоже квадратичная:
\[
  V^*(\vect{x}) = \vect{x}\T \mat{P}\,\vect{x},
\]
где $\mat{P}$ --- симметричная положительно определённая матрица.
\end{ideabox}

Подставляем в уравнение Беллмана:

\[
  \vect{x}\T \mat{P}\,\vect{x}
  = \min_{\vect{u}}
    \Big\{\vect{x}\T \mat{Q}\,\vect{x} + \vect{u}\T \mat{R}\,\vect{u}
          + (\mat{A}_d\vect{x} + \mat{B}_d\vect{u})\T \mat{P}\,
            (\mat{A}_d\vect{x} + \mat{B}_d\vect{u})\Big\}.
\]

Минимум по $\vect{u}$ ищется аналитически: производная по $\vect{u}$,
равная нулю.

\subsection{Шаг 1: найти оптимальное \texorpdfstring{$\vect{u}^*$}{u*}}

\begin{align*}
  \frac{\partial}{\partial \vect{u}}\Big[
    \vect{u}\T \mat{R}\vect{u} + (\mat{A}_d\vect{x} + \mat{B}_d\vect{u})\T
    \mat{P}(\mat{A}_d\vect{x} + \mat{B}_d\vect{u})
  \Big]
  &= 2\mat{R}\vect{u} + 2\mat{B}_d\T \mat{P}\,(\mat{A}_d\vect{x} + \mat{B}_d\vect{u})\\
  &= 0.
\end{align*}

Раскрываем:
\[
  (\mat{R} + \mat{B}_d\T \mat{P}\mat{B}_d)\,\vect{u} =
    -\mat{B}_d\T \mat{P}\mat{A}_d\,\vect{x}.
\]

\begin{formulabox}[title={Оптимальное управление через $\mat{P}$}]
\begin{equation}
  \vect{u}^* = -\mat{K}\,\vect{x},
  \qquad
  \mat{K} = (\mat{R} + \mat{B}_d\T \mat{P}\mat{B}_d)^{-1} \mat{B}_d\T \mat{P}\mat{A}_d.
  \label{eq:lqr_gain}
\end{equation}
\end{formulabox}

Это \textbf{линейный} закон: оптимальное управление --- линейная функция
текущего состояния. Чудо ЛКР именно в этом.

\subsection{Шаг 2: уравнение для $\mat{P}$ (DARE)}

Подставим $\vect{u}^* = -\mat{K}\vect{x}$ обратно в уравнение Беллмана
и упростим. После выкладок получим:

\begin{formulabox}[title={DARE --- Discrete Algebraic Riccati Equation}]
\begin{equation}
  \mat{P} \;=\; \mat{A}_d\T \mat{P}\mat{A}_d
              - \mat{A}_d\T \mat{P}\mat{B}_d
                (\mat{R} + \mat{B}_d\T \mat{P}\mat{B}_d)^{-1}
                \mat{B}_d\T \mat{P}\mat{A}_d
              + \mat{Q}.
  \label{eq:dare}
\end{equation}
\end{formulabox}

Это \emph{нелинейное матричное уравнение} относительно $\mat{P}$.

\section{Как решают DARE численно}

\begin{notebox}
Аналитически решить DARE невозможно (кроме скалярных случаев). На практике
используют численные методы:

\begin{itemize}
  \item \textbf{Итерации Риккати} --- запускаем рекурсию
        $\mat{P}_{k+1} = \text{RHS DARE}(\mat{P}_k)$ от любой начальной
        $\mat{P}_0$ (часто $\mat{Q}$). Сходится к решению, если объект
        управляем и $\mat{Q}, \mat{R}$ положительны.
  \item \textbf{Метод Шура} (Шурова декомпозиция гамильтоновой матрицы) ---
        стандарт в MATLAB и SciPy.
\end{itemize}

В нашем проекте --- наивные итерации (см. \filepath{pi\_nodes/control/\_linalg.py}),
fallback на \texttt{scipy.linalg.solve\_discrete\_are}, когда scipy есть.
\end{notebox}

\begin{codebox}[title={\filepath{pi\_nodes/control/\_linalg.py}, функция \texttt{solve\_discrete\_are}}]
\begin{lstlisting}[language=Python]
def solve_discrete_are(A, B, Q, R, max_iter=10_000, tol=1e-9):
    """Решает DARE итерациями Беллмана."""
    P = Q.copy()
    for _ in range(max_iter):
        BtP = B.T @ P
        K = np.linalg.solve(R + BtP @ B, BtP @ A)
        P_new = A.T @ P @ A - A.T @ P @ B @ K + Q
        if np.max(np.abs(P_new - P)) < tol:
            return P_new
        P = P_new
    raise RuntimeError("DARE did not converge")
\end{lstlisting}
\end{codebox}

После того как DARE решено и $\mat{P}$ найдена, считаем
$\mat{K}$ по \eqref{eq:lqr_gain}. Готово --- регулятор готов к работе.

\section{Какие $\mat{Q}, \mat{R}$ выбрать}

\begin{ideabox}
Выбор $\mat{Q}, \mat{R}$ --- это \emph{инженерное искусство}. ЛКР не
говорит, какие они должны быть; он только обещает \emph{оптимальность}
для тех, которые мы зададим.
\end{ideabox}

Практические рецепты:

\subsection{Правило Брайсона}

Для каждой компоненты $i$ положить $q_i = 1/x_{i,\max}^2$, где
$x_{i,\max}$ --- максимально \emph{допустимое} отклонение. Аналогично
для $r_j = 1/u_{j,\max}^2$. Это нормализует масштабы --- регулятор не
будет «думать», что отклонение в радианах ($\sim 0.1$) важнее, чем в
метрах ($\sim 5$).

\subsection{Дефолты проекта}

\begin{codebox}[title={\filepath{config.yaml}, секция \texttt{control.weights}}]
\begin{verbatim}
control:
  weights:
    Q_diag: [10.0, 10.0, 5.0, 1.0, 1.0]
    R_diag: [1.0, 1.0]
\end{verbatim}
\end{codebox}

\begin{tabular}{lll}
\toprule
Вес & Компонента & Смысл\\
\midrule
$10$ & $p_x$       & сильный штраф за «недокат» / «перекат»\\
$10$ & $p_y$       & сильный штраф за «съезд вбок»\\
$5$  & $\theta$    & средний штраф за кривой курс\\
$1$  & $v$         & лёгкий штраф за «не та скорость»\\
$1$  & $\omega$    & лёгкий штраф за «не то вращение»\\
$1$  & $u_v$       & вес на расход линейного управления\\
$1$  & $u_\omega$  & вес на расход углового управления\\
\bottomrule
\end{tabular}

\subsection{Что меняется, если\dots}

\begin{whatifbox}
\textbf{Q $\gg$ R} --- регулятор готов «топить педаль в пол», лишь бы
$\vect{x}$ был у цели. Быстро, но энергоёмко и с перерегулированием.

\textbf{Q $\ll$ R} --- регулятор «бережёт моторы», медленно подходит
к цели. Гладко, но долго.

\textbf{$\mat{Q}$ не-диагональная} --- разрешает «обмен» между
компонентами: можно сильно штрафовать, например, разницу $p_x - s_\text{ref}$,
а не каждую компоненту отдельно. Усложняет интерпретацию, редко нужно.

\textbf{$q_\theta = 0$} --- регулятору всё равно, куда смотрит робот.
Странно, но математически валидно: ЛКР решит DARE и выдаст $\mat{K}$,
где курс будет «болтаться» свободно.

\textbf{$r_i = 0$} --- неуправляемо, DARE не сходится, $\mat{P}$
вырождается. \emph{Не делать так.}
\end{whatifbox}

\section{Замкнутая система}

\begin{ideabox}
Подставим закон $\vect{u}_k = -\mat{K}\vect{x}_k$ в исходную модель:
\[
  \vect{x}_{k+1} = \mat{A}_d\vect{x}_k + \mat{B}_d(-\mat{K}\vect{x}_k)
                 = (\mat{A}_d - \mat{B}_d\mat{K})\,\vect{x}_k.
\]
Это \textbf{замкнутая система}. Её собственные значения $\lambda(\mat{A}_d
- \mat{B}_d \mat{K})$ --- \emph{полюса} замкнутого контура. ЛКР гарантирует,
что все они лежат внутри единичного круга ($|\lambda| < 1$) --- замкнутая
система \emph{устойчива}.
\end{ideabox}

Проверка в коде:

\begin{codebox}[title={\filepath{pi\_nodes/control/mpc\_controller.py:332}}]
\begin{lstlisting}[language=Python]
def closed_loop_eigenvalues(self) -> np.ndarray:
    """|λ| < 1 ⇔ closed-loop x_{k+1} = (Ad - Bd K) x is stable."""
    return np.linalg.eigvals(self.Ad - self.Bd @ self.K_first)

def is_stable(self) -> bool:
    return bool(np.all(np.abs(self.closed_loop_eigenvalues()) < 1.0 - 1e-9))
\end{lstlisting}
\end{codebox}

\section{Слабое место: ограничения}

ЛКР \textbf{не знает} о том, что $u_v \le 0.30$~м/с и $|u_\omega| \le
2$~рад/с. Если задать $\vect{x}_0$ далеко от цели, закон $-\mat{K}\vect{x}_0$
выдаст $\vect{u}$, который физически невозможен.

\begin{warnbox}
В коде \texttt{LQRController.step()} это решают \emph{пост-обработкой}:
после $-\mat{K}\vect{x}$ результат прогоняется через
\texttt{np.clip(u, u\_min, u\_max)}. Это работает, но \emph{не оптимально}
--- ЛКР проектировал под полную свободу. Когда $\vect{u}$ упирается в
насыщение, его поведение перестаёт быть оптимальным.

Это --- \emph{мотивация для MPC}. MPC формулирует ограничения внутри
оптимизации, а не накладывает их сверху.
\end{warnbox}

\section{Что мы получили}

\begin{itemize}
  \item \textbf{ЛКР} --- линейный оптимальный регулятор:
        $\vect{u}_k = -\mat{K}\vect{x}_k$.
  \item Матрица $\mat{K}$ считается \emph{оффлайн} через DARE
        \eqref{eq:dare} --- одно матричное уравнение, решается итерациями
        либо Шуром.
  \item ЛКР гарантирует устойчивость замкнутой системы (все полюса
        внутри единичного круга).
  \item Выбор $\mat{Q}, \mat{R}$ --- инженерное искусство: правило
        Брайсона, диагональные веса.
  \item ЛКР \textbf{не знает} о box-ограничениях на $\vect{u}$. Клип
        сверху --- костыль, мотивация для MPC.
\end{itemize}

Идём дальше: MPC.
